\documentclass[a4paper,11pt]{article}
\usepackage{amssymb, enumitem}
\setlist[enumerate,1]{label={\bfseries\arabic*.},ref={\bfseries{Ejercicio \arabic*}}}
\setlist[enumerate,2]{label={\bfseries(\alph*)},ref={\bfseries(\alph*)}}

\usepackage{tikz}

\parindent 0cm
\usepackage{amssymb,amsmath,amsthm,latexsym,epsfig,euscript,multicol}
\usepackage{graphbox}

\usepackage[utf8x]{inputenc}
\usepackage{listings,xcolor,bm}


\definecolor{mygreen}{rgb}{0,0.6,0}
\definecolor{mygray}{rgb}{0.5,0.5,0.5}
\definecolor{mymauve}{rgb}{0.58,0,0.82}
\lstset{
  backgroundcolor=\color{white},   % choose the background color; you must add
  basicstyle=\ttfamily,      % the size of the fonts that are used for the code
  breakatwhitespace=true,         % sets if automatic breaks should only happen at whitespace
  breaklines=true,                 % sets automatic line breaking
  captionpos=b,                    % sets the caption-position to bottom
  commentstyle=\color{mygreen},    % comment style
  deletekeywords={...},            % if you want to delete keywords from the given language
  escapeinside={\%*}{*)},          % if you want to add LaTeX within your code
  extendedchars=true,              % lets you use non-ASCII characters; for 8-bits encodings only, does not work with UTF-8
  firstnumber=1,                % start line enumeration with line 1000
  frame=single,	                   % adds a frame around the code
  keepspaces=true,                 % keeps spaces in text, useful for keeping indentation of code (possibly needs columns=flexible)
  keywordstyle=\color{blue},       % keyword style
  language=Python,                 % the language of the code
  morekeywords={*,...},            % if you want to add more keywords to the set
  numbers=none,                    % where to put the line-numbers; possible values are (none, left, right)
  numbersep=5pt,                   % how far the line-numbers are from the code
  numberstyle=\tiny\color{mygray}, % the style that is used for the line-numbers
  rulecolor=\color{black},         % if not set, the frame-color may be changed on line-breaks within not-black text (e.g. comments (green here))
  showspaces=false,                % show spaces everywhere adding particular underscores; it overrides 'showstringspaces'
  showstringspaces=false,          % underline spaces within strings only
  showtabs=false,                  % show tabs within strings adding particular underscores
  stepnumber=5,                    % the step between two line-numbers. If it's 1, each line will be numbered
  stringstyle=\color{mymauve},     % string literal style
  tabsize=4,	                   % sets default tabsize to 2 spaces
  title=\lstname,                  % show the filename of files included with \lstinputlisting; also try caption instead of title
  belowskip=0em
}
% Caracteres especiales
\def\A{\mathbb{A}}
\def\C{\mathbb{C}}
\def \N{\mathbb{N}}
\def \P{\mathbb{P}}
\def \Q{\mathbb{Q}}
\def \R{\mathbb{R}}
\def \Z{\mathbb{Z}}
\def \sen{\textrm{sen}}

\def\Np{$\N$}
\def\Zp{$\Z$}
\def\Qp{$\Q$}
\def\Rp{$\R$}
\def\Cp{$\C$}

\def\bb{\bm{b}}
\def\bu{\bm{u}}
\def\bv{\bm{v}}
\def\bx{\bm{x}}
\def\bA{\bm{A}}
\def\bB{\bm{B}}
\def\bD{\bm{D}}
\def\bE{\bm{E}}
\def\bM{\bm{M}}
\def\bT{\bm{T}}


\def\K{\textrm{K}}
\def\V{\textrm{V}}
\def\S{\textrm{S}}

\def\degres{$^\circ$}

\newcount\todno
\def\no{\global\advance\todno by 1 \the\todno}

\topmargin-2cm \vsize 29.5cm \hsize 21cm
\setlength{\textwidth}{16.75cm}\setlength{\textheight}{23.5cm}
\setlength{\oddsidemargin}{0.0cm}
\setlength{\evensidemargin}{0.0cm}


\theoremstyle{definition}
\newtheorem{ejer}{Ejercicio}
\newcommand{\bej}{\begin{ejer}}
\newcommand{\fej}{\end{ejer}}

\begin{document}

\centerline{{\small Universidad de Buenos Aires - Facultad de Ciencias Exactas y Naturales - Ciencias de Datos}}

\vskip 0.2cm

\hrule

\vskip 0.2cm

 \centerline{{\bf\Large{\sc Laboratorio de Datos}}}

 \vskip 0.2cm

 \centerline{\ttfamily Primer Cuatrimestre 2024}

\vskip 0.2cm

 \hrule

 \bigskip
 \centerline{\bf Práctica N$^\circ$ 7: Clustering.}
 \bigskip

%-----------------------------------------------------------
\begin{enumerate}
\item Dada la siguiente tabla de datos
\begin{table}[h]
\begin{center}
\begin{tabular}{|c|c|c|c|c|c|c|c|}
\hline
x & -1 & 0 & 1 & 8/5 & 2 & 3 & 4  \\
y  & 2 & 1& 2& 2& 1& 2 & 1 \\
\hline
\end{tabular}
\end{center}
\end{table}

Utilizar ``a mano'' el m\'etodo de $k$-medias para agrupar los datos en 2 clusters.
\begin{enumerate}
\item Comenzando con $b_1 =(1,2)  $ y $b_2=(3,2)  $
\item Comenzando con $b_1=(0,1)   $ y $b_2= (3,2) $
\end{enumerate}

?`Se obtiene la misma clasificaci\'on? ?`Alguna de las clasificaciones obtenidas le parece m\'as apropiada?

\item Considere los datasetets \texttt{p7-data1.csv} y \texttt{p7-data2.csv} de datos artificialmente generados.

\begin{enumerate}
	\item Abra cada dataset en Python y genere un diagrama de dispersi\'on (scatter plot) para cada uno.
	\item Analizando los gr\'aficos ``a mano'' considere cu\'antos clusters est\'an presentes.
	\item Pruebe ejecutar el comando \texttt{KMeans} con la cantidad de clusters que detect\'o. Analizar el comportamiento del procedimiento en cada caso.
	
\end{enumerate}

%-----------------------------------------------------------


\item Considerar el dataset \texttt{p7-iris.txt} (para leer el archivo, observar que los datos est\'an separados por tabulaciones). En este ejercicio trataremos de identificar las distintas subespecies.
\begin{enumerate}
	\item Cargue el archivo \texttt{p7-iris.txt}.
	\item Grafique en un diagrama de dispersi\'on la longitud del p\'etalo vs el ancho del p\'etalo.
	\item Efect\'ue un \textit{clustering} k-medias con el comando \texttt{KMeans} de los datos basados en las cuatro columnas de datos, considere $k=3$ clusters.
	\item Repita el inciso b) coloreando en funci\'on del \'indice de cluster obtenido.
	\item Eval\'ue el error de clustering en funci\'on de la siguiente f\'ormula (within-cluster sum of squares, WCSS):
	$$
	WCSS = \sum_{i=}^k \sum_{x \in C_i} \| x - \mu_i\|^2
	$$
	donde $C_i$ representa el cluster i-\'esimo y $\mu_i$ es el centroide de dicho cluster, definido como
	$$
	\mu_i = \frac{1}{\#C_i}\sum_{x \in C_i} x.
	$$
	%Octave ofrece una forma de calcular esto de forma m\'as directa, lea la documentaci\'on del comando \textbf{KMeans}
	\textit{Python ofrece una forma de calcular esto de forma directa.}
	
	 \textit{(Mirar el archivo \texttt{p7-ejercicioPetalos.ipynb}.)}
	\item Repita el ensayo para distintos valores de $k$, entre 1 y 10, graficando el $WCSS$ para cada valor de $k$. Analizar el mejor valor de $k$ posible teniendo en cuenta
	un compromiso entre ``complejidad'' (es decir, cantidad de clusters) y nivel de error (es decir, el WCSS).
	
\end{enumerate}

%-----------------------------------------------------------



\item Consideremos el dataset de datos artificiales \texttt{p7-dataSinEscalar.csv}.
\begin{enumerate}
	\item Cargar los datos  y graf\'icarlos.
	\item A priori y mirando el gr\'afico, determine la cantidad de clusters que puede detectar en los mismos e imagine inicialmente c\'omo debieran ser esos clusters.
	\item Realizar un clustering k-medias con el valor de $k$ antes determinado.
	\item ?`Considera satisfactorio el clustering obtenido? ?`Representa lo que usted esperaba?
	\item Uno de los problemas que tenemos es que el m\'etodo de k-medias es muy sensible a las diferencias de escala entre las dimensiones. Una forma de corregir eso es re-escalando las variables de
	forma tal que todas se muevan en el mismo rango. Por ejemplo, podemos conseguir eso efectuando una normalizaci\'on como sigue:
	
	$$
	X_{ij} = \frac{X_{ij} - min(X_{\cdot j}) }{max(X_{\cdot j}) - min(X_{\cdot j})}.
	$$
	
	De esta manera, logramos que los datos de cada columna caigan entre 0 y 1. Normalice los datos siguiendo este criterio.
	
	\textit{(Mirar en Python el comando \texttt{MinMaxScaler})}
	\item Vuelva a correr el procedimiento de clustering, tome las etiquetas de clustering obtenidos y grafique los datos originales con un color que dependa del clustering obtenido con los datos escalados.
	
\end{enumerate}

%-----------------------------------------------------------

\item (opcional) Implementar el algoritmo DBSCAN para analizar los sets de datos anteriores. Comparar los resultados con los obtenidos usando k-medias.

%-----------------------------------------------------------




\end{enumerate}

\end{document}


\item Considerando los datos de la tabla
\begin{enumerate}
	\item Cargar los siguientes datos en Python y graficarlos mediante un diagrama de dispersi\'on. Procurar que cada punto aparezca etiquetado con su \'indice.
	\begin{center}
		\begin{tabular}{lllllllllll}
			\'indice & 1 & 2  & 3  & 4  & 5  & 6  & 7  & 8  & 9  & 10 \\ \hline
			x        & 5 & 10 & 15 & 24 & 30 & 85 & 71 & 60 & 70 & 80\\
			y        & 3 & 15 & 12 & 10 & 30 & 70 & 80 & 78 & 55 & 91  \\
		\end{tabular}
	\end{center}
	\item ``A mano'' e informalmente, identificar agrupaciones jer\'arquicas en el gr\'afico.
	\item Mediante un clustering jer\'arquico, basado en la distancia eucl\'idea y usando como criterio de agregaci\'on el esquema ``single'', graficar un dendrograma de los datos.
	\item Marcar los puntos que tomar\'ian en caso de querer considerar dos clusters.
	\item \'idem inciso anterior pero ahora con 4 clusters.
\end{enumerate}


%%-----------------------------------------------------------

\item Utilice clustering jer\'arquico para poder obtener un clustering del dataset \texttt{p7-data1.csv} con dos clusters, compararlo con aquel obtenido usando k-medias. Hacer lo mismo para el dataset \texttt{p7-data2.csv}

%%-----------------------------------------------------------

